% paper67-code-review-automation.tex
% Automating Code Review with Judgment-Based Analysis
% JuGeo Paper Series — Paper 67

\documentclass[11pt]{article}
% jugeo-common.tex — shared preamble for all Judgment Geometry papers
% Include via: % jugeo-common.tex — shared preamble for all Judgment Geometry papers
% Include via: % jugeo-common.tex — shared preamble for all Judgment Geometry papers
% Include via: \input{jugeo-common}

% ── Packages ────────────────────────────────────────────────────────────
\usepackage{amsmath,amssymb,amsthm,mathtools}
\usepackage{tikz-cd}
\usepackage{listings}
\usepackage{xcolor}
\usepackage{xspace}
\usepackage{booktabs}
\usepackage{multirow}
\usepackage{enumitem}
\usepackage[expansion=false]{microtype}
\usepackage{hyperref}
\usepackage{cleveref}
% \usepackage{thmtools}  % disabled: not available in this TeX installation
\usepackage{float}
\usepackage{caption}
\usepackage{subcaption}
\usepackage{tabularx}
\usepackage{stmaryrd}       % \llbracket, \rrbracket
\usepackage{bbm}            % \mathbbm{1}

% ── Page geometry ───────────────────────────────────────────────────────
\usepackage[margin=1in]{geometry}

% ── Theorem environments ────────────────────────────────────────────────
\theoremstyle{plain}
\newtheorem{theorem}{Theorem}[section]
\newtheorem{lemma}[theorem]{Lemma}
\newtheorem{proposition}[theorem]{Proposition}
\newtheorem{corollary}[theorem]{Corollary}

\theoremstyle{definition}
\newtheorem{definition}[theorem]{Definition}
\newtheorem{example}[theorem]{Example}
\newtheorem{construction}[theorem]{Construction}

\theoremstyle{remark}
\newtheorem{remark}[theorem]{Remark}
\newtheorem{notation}[theorem]{Notation}

% ── Python listings style ──────────────────────────────────────────────
\definecolor{jugeoBlue}{HTML}{2563EB}
\definecolor{jugeoGreen}{HTML}{059669}
\definecolor{jugeoRed}{HTML}{DC2626}
\definecolor{jugeoPurple}{HTML}{7C3AED}
\definecolor{jugeoGray}{HTML}{6B7280}
\definecolor{jugeoBg}{HTML}{F8FAFC}

\lstdefinestyle{jugeo-python}{
  language=Python,
  basicstyle=\ttfamily\small,
  keywordstyle=\color{jugeoBlue}\bfseries,
  stringstyle=\color{jugeoGreen},
  commentstyle=\color{jugeoGray}\itshape,
  numberstyle=\tiny\color{jugeoGray},
  backgroundcolor=\color{jugeoBg},
  numbers=left,
  numbersep=6pt,
  frame=single,
  framerule=0.4pt,
  rulecolor=\color{jugeoGray!40},
  breaklines=true,
  breakatwhitespace=true,
  showstringspaces=false,
  tabsize=4,
  xleftmargin=1.5em,
  framexleftmargin=1.5em,
  morekeywords={dataclass,frozen,field,Enum,IntEnum,auto,Any,
                Mapping,Sequence,Optional,match,case},
  literate={→}{{$\to$}}1 {←}{{$\leftarrow$}}1
           {≤}{{$\leq$}}1 {≥}{{$\geq$}}1 {≠}{{$\neq$}}1
           {∈}{{$\in$}}1 {∀}{{$\forall$}}1 {∃}{{$\exists$}}1
           {⊕}{{$\oplus$}}1 {⊖}{{$\ominus$}}1,
  escapeinside={(*@}{@*)},
}
\lstset{style=jugeo-python}

\lstdefinestyle{jugeo-output}{
  basicstyle=\ttfamily\small,
  backgroundcolor=\color{jugeoBg},
  frame=single,
  framerule=0.4pt,
  rulecolor=\color{jugeoGray!40},
  breaklines=true,
  showstringspaces=false,
  xleftmargin=1.5em,
  framexleftmargin=0.5em,
  numbers=none,
}

% ── JuGeo-specific macros ──────────────────────────────────────────────

% Judgment 8-tuple
\newcommand{\J}{\mathcal{J}}
\newcommand{\Jud}[8]{(#1,\,#2,\,#3,\,#4,\,#5,\,#6,\,#7,\,#8)}
\newcommand{\JudShort}[1]{\J_{#1}}

% Components
\newcommand{\coord}{\mathsf{c}}            % coordinate
\newcommand{\prop}{\varphi}                 % proposition
\newcommand{\carrier}{\mathcal{A}}          % carrier type
\newcommand{\evid}{\mathcal{E}}             % evidence bundle
\newcommand{\oblig}{\mathcal{O}}            % obligations
\newcommand{\obstr}{\mathcal{B}}            % obstructions
\newcommand{\trust}{\mathcal{T}}            % trust annotation
\newcommand{\prov}{\Pi}                     % provenance

% Trust algebra
\newcommand{\Talg}{\mathfrak{T}}            % trust ordered algebra
\newcommand{\Eadm}{\mathcal{E}_{\mathrm{adm}}}
\newcommand{\tleq}{\preceq}                 % trust ordering
\newcommand{\tjoin}{\oplus}                 % conservative join
\newcommand{\tatten}{\ominus}               % attenuation
\newcommand{\tpromote}{\uparrow_\pi}        % promotion
\newcommand{\tdemote}{\downarrow_\chi}       % demotion

% Trust tiers
\newcommand{\Tcontra}{\ensuremath{\mathsf{CONTRADICTED}}}
\newcommand{\Tunver}{\ensuremath{\mathsf{UNVERIFIED}}}
\newcommand{\Tcopilot}{\ensuremath{\mathsf{COPILOT}}}
\newcommand{\Toracle}{\ensuremath{\mathsf{ORACLE}}}
\newcommand{\Truntime}{\ensuremath{\mathsf{RUNTIME}}}
\newcommand{\Tsolver}{\ensuremath{\mathsf{SOLVER}}}
\newcommand{\Tproof}{\ensuremath{\mathsf{PROOF}}}

% Site and geometry
\newcommand{\Site}{\mathcal{S}}             % semantic site
\newcommand{\Cov}{\mathrm{Cov}}            % covering family
\newcommand{\Ob}{\mathrm{Ob}}              % objects
\newcommand{\Mor}{\mathrm{Mor}}            % morphisms
\newcommand{\res}{\mathrm{res}}            % restriction
\newcommand{\Cech}{\check{C}}              % Čech complex
\newcommand{\Hcoh}[1]{H^{#1}}             % cohomology group

% Descent
\newcommand{\descent}{\mathrm{desc}}
\newcommand{\glue}{\mathrm{glue}}
\newcommand{\overlap}[2]{#1 \cap #2}

% Semantic moves
\newcommand{\VERIFY}{\textsc{Verify}}
\newcommand{\CONSTRUCT}{\textsc{Construct}}
\newcommand{\NORMALIZE}{\textsc{Normalize}}
\newcommand{\GLUE}{\textsc{Glue}}
\newcommand{\RESTRICT}{\textsc{Restrict}}
\newcommand{\TRANSPORT}{\textsc{Transport}}
\newcommand{\REPAIR}{\textsc{Repair}}
\newcommand{\PROMOTE}{\textsc{Promote}}
\newcommand{\CHALLENGE}{\textsc{Challenge}}
\newcommand{\ESCALATE}{\textsc{Escalate}}

% Fragments
\newcommand{\QFLIA}{\texttt{QF\_LIA}}
\newcommand{\QFLRA}{\texttt{QF\_LRA}}
\newcommand{\QFBV}{\texttt{QF\_BV}}
\newcommand{\QFUF}{\texttt{QF\_UF}}

% Misc
\newcommand{\Python}{\textsc{Python}}
\newcommand{\JuGeo}{\textsc{JuGeo}}
\newcommand{\LEAN}{\textsc{Lean}\,4}
\newcommand{\Fstar}{F$^{\star}$}
\newcommand{\Coq}{\textsc{Coq}}
\newcommand{\Dafny}{\textsc{Dafny}}

% Common shortcuts
\newcommand{\ie}{i.e.\@\xspace}
\newcommand{\eg}{e.g.\@\xspace}
\newcommand{\cf}{cf.\@\xspace}
\newcommand{\wrt}{w.r.t.\@\xspace}

% Evaluation shortcuts
\newcommand{\numcases}{300}
\newcommand{\numspec}{100}
\newcommand{\numeq}{100}
\newcommand{\numbug}{100}
\newcommand{\medtime}{6.5\,ms}
\newcommand{\totaltime}{1.949\,s}


% ── Packages ────────────────────────────────────────────────────────────
\usepackage{amsmath,amssymb,amsthm,mathtools}
\usepackage{tikz-cd}
\usepackage{listings}
\usepackage{xcolor}
\usepackage{xspace}
\usepackage{booktabs}
\usepackage{multirow}
\usepackage{enumitem}
\usepackage[expansion=false]{microtype}
\usepackage{hyperref}
\usepackage{cleveref}
% \usepackage{thmtools}  % disabled: not available in this TeX installation
\usepackage{float}
\usepackage{caption}
\usepackage{subcaption}
\usepackage{tabularx}
\usepackage{stmaryrd}       % \llbracket, \rrbracket
\usepackage{bbm}            % \mathbbm{1}

% ── Page geometry ───────────────────────────────────────────────────────
\usepackage[margin=1in]{geometry}

% ── Theorem environments ────────────────────────────────────────────────
\theoremstyle{plain}
\newtheorem{theorem}{Theorem}[section]
\newtheorem{lemma}[theorem]{Lemma}
\newtheorem{proposition}[theorem]{Proposition}
\newtheorem{corollary}[theorem]{Corollary}

\theoremstyle{definition}
\newtheorem{definition}[theorem]{Definition}
\newtheorem{example}[theorem]{Example}
\newtheorem{construction}[theorem]{Construction}

\theoremstyle{remark}
\newtheorem{remark}[theorem]{Remark}
\newtheorem{notation}[theorem]{Notation}

% ── Python listings style ──────────────────────────────────────────────
\definecolor{jugeoBlue}{HTML}{2563EB}
\definecolor{jugeoGreen}{HTML}{059669}
\definecolor{jugeoRed}{HTML}{DC2626}
\definecolor{jugeoPurple}{HTML}{7C3AED}
\definecolor{jugeoGray}{HTML}{6B7280}
\definecolor{jugeoBg}{HTML}{F8FAFC}

\lstdefinestyle{jugeo-python}{
  language=Python,
  basicstyle=\ttfamily\small,
  keywordstyle=\color{jugeoBlue}\bfseries,
  stringstyle=\color{jugeoGreen},
  commentstyle=\color{jugeoGray}\itshape,
  numberstyle=\tiny\color{jugeoGray},
  backgroundcolor=\color{jugeoBg},
  numbers=left,
  numbersep=6pt,
  frame=single,
  framerule=0.4pt,
  rulecolor=\color{jugeoGray!40},
  breaklines=true,
  breakatwhitespace=true,
  showstringspaces=false,
  tabsize=4,
  xleftmargin=1.5em,
  framexleftmargin=1.5em,
  morekeywords={dataclass,frozen,field,Enum,IntEnum,auto,Any,
                Mapping,Sequence,Optional,match,case},
  literate={→}{{$\to$}}1 {←}{{$\leftarrow$}}1
           {≤}{{$\leq$}}1 {≥}{{$\geq$}}1 {≠}{{$\neq$}}1
           {∈}{{$\in$}}1 {∀}{{$\forall$}}1 {∃}{{$\exists$}}1
           {⊕}{{$\oplus$}}1 {⊖}{{$\ominus$}}1,
  escapeinside={(*@}{@*)},
}
\lstset{style=jugeo-python}

\lstdefinestyle{jugeo-output}{
  basicstyle=\ttfamily\small,
  backgroundcolor=\color{jugeoBg},
  frame=single,
  framerule=0.4pt,
  rulecolor=\color{jugeoGray!40},
  breaklines=true,
  showstringspaces=false,
  xleftmargin=1.5em,
  framexleftmargin=0.5em,
  numbers=none,
}

% ── JuGeo-specific macros ──────────────────────────────────────────────

% Judgment 8-tuple
\newcommand{\J}{\mathcal{J}}
\newcommand{\Jud}[8]{(#1,\,#2,\,#3,\,#4,\,#5,\,#6,\,#7,\,#8)}
\newcommand{\JudShort}[1]{\J_{#1}}

% Components
\newcommand{\coord}{\mathsf{c}}            % coordinate
\newcommand{\prop}{\varphi}                 % proposition
\newcommand{\carrier}{\mathcal{A}}          % carrier type
\newcommand{\evid}{\mathcal{E}}             % evidence bundle
\newcommand{\oblig}{\mathcal{O}}            % obligations
\newcommand{\obstr}{\mathcal{B}}            % obstructions
\newcommand{\trust}{\mathcal{T}}            % trust annotation
\newcommand{\prov}{\Pi}                     % provenance

% Trust algebra
\newcommand{\Talg}{\mathfrak{T}}            % trust ordered algebra
\newcommand{\Eadm}{\mathcal{E}_{\mathrm{adm}}}
\newcommand{\tleq}{\preceq}                 % trust ordering
\newcommand{\tjoin}{\oplus}                 % conservative join
\newcommand{\tatten}{\ominus}               % attenuation
\newcommand{\tpromote}{\uparrow_\pi}        % promotion
\newcommand{\tdemote}{\downarrow_\chi}       % demotion

% Trust tiers
\newcommand{\Tcontra}{\ensuremath{\mathsf{CONTRADICTED}}}
\newcommand{\Tunver}{\ensuremath{\mathsf{UNVERIFIED}}}
\newcommand{\Tcopilot}{\ensuremath{\mathsf{COPILOT}}}
\newcommand{\Toracle}{\ensuremath{\mathsf{ORACLE}}}
\newcommand{\Truntime}{\ensuremath{\mathsf{RUNTIME}}}
\newcommand{\Tsolver}{\ensuremath{\mathsf{SOLVER}}}
\newcommand{\Tproof}{\ensuremath{\mathsf{PROOF}}}

% Site and geometry
\newcommand{\Site}{\mathcal{S}}             % semantic site
\newcommand{\Cov}{\mathrm{Cov}}            % covering family
\newcommand{\Ob}{\mathrm{Ob}}              % objects
\newcommand{\Mor}{\mathrm{Mor}}            % morphisms
\newcommand{\res}{\mathrm{res}}            % restriction
\newcommand{\Cech}{\check{C}}              % Čech complex
\newcommand{\Hcoh}[1]{H^{#1}}             % cohomology group

% Descent
\newcommand{\descent}{\mathrm{desc}}
\newcommand{\glue}{\mathrm{glue}}
\newcommand{\overlap}[2]{#1 \cap #2}

% Semantic moves
\newcommand{\VERIFY}{\textsc{Verify}}
\newcommand{\CONSTRUCT}{\textsc{Construct}}
\newcommand{\NORMALIZE}{\textsc{Normalize}}
\newcommand{\GLUE}{\textsc{Glue}}
\newcommand{\RESTRICT}{\textsc{Restrict}}
\newcommand{\TRANSPORT}{\textsc{Transport}}
\newcommand{\REPAIR}{\textsc{Repair}}
\newcommand{\PROMOTE}{\textsc{Promote}}
\newcommand{\CHALLENGE}{\textsc{Challenge}}
\newcommand{\ESCALATE}{\textsc{Escalate}}

% Fragments
\newcommand{\QFLIA}{\texttt{QF\_LIA}}
\newcommand{\QFLRA}{\texttt{QF\_LRA}}
\newcommand{\QFBV}{\texttt{QF\_BV}}
\newcommand{\QFUF}{\texttt{QF\_UF}}

% Misc
\newcommand{\Python}{\textsc{Python}}
\newcommand{\JuGeo}{\textsc{JuGeo}}
\newcommand{\LEAN}{\textsc{Lean}\,4}
\newcommand{\Fstar}{F$^{\star}$}
\newcommand{\Coq}{\textsc{Coq}}
\newcommand{\Dafny}{\textsc{Dafny}}

% Common shortcuts
\newcommand{\ie}{i.e.\@\xspace}
\newcommand{\eg}{e.g.\@\xspace}
\newcommand{\cf}{cf.\@\xspace}
\newcommand{\wrt}{w.r.t.\@\xspace}

% Evaluation shortcuts
\newcommand{\numcases}{300}
\newcommand{\numspec}{100}
\newcommand{\numeq}{100}
\newcommand{\numbug}{100}
\newcommand{\medtime}{6.5\,ms}
\newcommand{\totaltime}{1.949\,s}


% ── Packages ────────────────────────────────────────────────────────────
\usepackage{amsmath,amssymb,amsthm,mathtools}
\usepackage{tikz-cd}
\usepackage{listings}
\usepackage{xcolor}
\usepackage{xspace}
\usepackage{booktabs}
\usepackage{multirow}
\usepackage{enumitem}
\usepackage[expansion=false]{microtype}
\usepackage{hyperref}
\usepackage{cleveref}
% \usepackage{thmtools}  % disabled: not available in this TeX installation
\usepackage{float}
\usepackage{caption}
\usepackage{subcaption}
\usepackage{tabularx}
\usepackage{stmaryrd}       % \llbracket, \rrbracket
\usepackage{bbm}            % \mathbbm{1}

% ── Page geometry ───────────────────────────────────────────────────────
\usepackage[margin=1in]{geometry}

% ── Theorem environments ────────────────────────────────────────────────
\theoremstyle{plain}
\newtheorem{theorem}{Theorem}[section]
\newtheorem{lemma}[theorem]{Lemma}
\newtheorem{proposition}[theorem]{Proposition}
\newtheorem{corollary}[theorem]{Corollary}

\theoremstyle{definition}
\newtheorem{definition}[theorem]{Definition}
\newtheorem{example}[theorem]{Example}
\newtheorem{construction}[theorem]{Construction}

\theoremstyle{remark}
\newtheorem{remark}[theorem]{Remark}
\newtheorem{notation}[theorem]{Notation}

% ── Python listings style ──────────────────────────────────────────────
\definecolor{jugeoBlue}{HTML}{2563EB}
\definecolor{jugeoGreen}{HTML}{059669}
\definecolor{jugeoRed}{HTML}{DC2626}
\definecolor{jugeoPurple}{HTML}{7C3AED}
\definecolor{jugeoGray}{HTML}{6B7280}
\definecolor{jugeoBg}{HTML}{F8FAFC}

\lstdefinestyle{jugeo-python}{
  language=Python,
  basicstyle=\ttfamily\small,
  keywordstyle=\color{jugeoBlue}\bfseries,
  stringstyle=\color{jugeoGreen},
  commentstyle=\color{jugeoGray}\itshape,
  numberstyle=\tiny\color{jugeoGray},
  backgroundcolor=\color{jugeoBg},
  numbers=left,
  numbersep=6pt,
  frame=single,
  framerule=0.4pt,
  rulecolor=\color{jugeoGray!40},
  breaklines=true,
  breakatwhitespace=true,
  showstringspaces=false,
  tabsize=4,
  xleftmargin=1.5em,
  framexleftmargin=1.5em,
  morekeywords={dataclass,frozen,field,Enum,IntEnum,auto,Any,
                Mapping,Sequence,Optional,match,case},
  literate={→}{{$\to$}}1 {←}{{$\leftarrow$}}1
           {≤}{{$\leq$}}1 {≥}{{$\geq$}}1 {≠}{{$\neq$}}1
           {∈}{{$\in$}}1 {∀}{{$\forall$}}1 {∃}{{$\exists$}}1
           {⊕}{{$\oplus$}}1 {⊖}{{$\ominus$}}1,
  escapeinside={(*@}{@*)},
}
\lstset{style=jugeo-python}

\lstdefinestyle{jugeo-output}{
  basicstyle=\ttfamily\small,
  backgroundcolor=\color{jugeoBg},
  frame=single,
  framerule=0.4pt,
  rulecolor=\color{jugeoGray!40},
  breaklines=true,
  showstringspaces=false,
  xleftmargin=1.5em,
  framexleftmargin=0.5em,
  numbers=none,
}

% ── JuGeo-specific macros ──────────────────────────────────────────────

% Judgment 8-tuple
\newcommand{\J}{\mathcal{J}}
\newcommand{\Jud}[8]{(#1,\,#2,\,#3,\,#4,\,#5,\,#6,\,#7,\,#8)}
\newcommand{\JudShort}[1]{\J_{#1}}

% Components
\newcommand{\coord}{\mathsf{c}}            % coordinate
\newcommand{\prop}{\varphi}                 % proposition
\newcommand{\carrier}{\mathcal{A}}          % carrier type
\newcommand{\evid}{\mathcal{E}}             % evidence bundle
\newcommand{\oblig}{\mathcal{O}}            % obligations
\newcommand{\obstr}{\mathcal{B}}            % obstructions
\newcommand{\trust}{\mathcal{T}}            % trust annotation
\newcommand{\prov}{\Pi}                     % provenance

% Trust algebra
\newcommand{\Talg}{\mathfrak{T}}            % trust ordered algebra
\newcommand{\Eadm}{\mathcal{E}_{\mathrm{adm}}}
\newcommand{\tleq}{\preceq}                 % trust ordering
\newcommand{\tjoin}{\oplus}                 % conservative join
\newcommand{\tatten}{\ominus}               % attenuation
\newcommand{\tpromote}{\uparrow_\pi}        % promotion
\newcommand{\tdemote}{\downarrow_\chi}       % demotion

% Trust tiers
\newcommand{\Tcontra}{\ensuremath{\mathsf{CONTRADICTED}}}
\newcommand{\Tunver}{\ensuremath{\mathsf{UNVERIFIED}}}
\newcommand{\Tcopilot}{\ensuremath{\mathsf{COPILOT}}}
\newcommand{\Toracle}{\ensuremath{\mathsf{ORACLE}}}
\newcommand{\Truntime}{\ensuremath{\mathsf{RUNTIME}}}
\newcommand{\Tsolver}{\ensuremath{\mathsf{SOLVER}}}
\newcommand{\Tproof}{\ensuremath{\mathsf{PROOF}}}

% Site and geometry
\newcommand{\Site}{\mathcal{S}}             % semantic site
\newcommand{\Cov}{\mathrm{Cov}}            % covering family
\newcommand{\Ob}{\mathrm{Ob}}              % objects
\newcommand{\Mor}{\mathrm{Mor}}            % morphisms
\newcommand{\res}{\mathrm{res}}            % restriction
\newcommand{\Cech}{\check{C}}              % Čech complex
\newcommand{\Hcoh}[1]{H^{#1}}             % cohomology group

% Descent
\newcommand{\descent}{\mathrm{desc}}
\newcommand{\glue}{\mathrm{glue}}
\newcommand{\overlap}[2]{#1 \cap #2}

% Semantic moves
\newcommand{\VERIFY}{\textsc{Verify}}
\newcommand{\CONSTRUCT}{\textsc{Construct}}
\newcommand{\NORMALIZE}{\textsc{Normalize}}
\newcommand{\GLUE}{\textsc{Glue}}
\newcommand{\RESTRICT}{\textsc{Restrict}}
\newcommand{\TRANSPORT}{\textsc{Transport}}
\newcommand{\REPAIR}{\textsc{Repair}}
\newcommand{\PROMOTE}{\textsc{Promote}}
\newcommand{\CHALLENGE}{\textsc{Challenge}}
\newcommand{\ESCALATE}{\textsc{Escalate}}

% Fragments
\newcommand{\QFLIA}{\texttt{QF\_LIA}}
\newcommand{\QFLRA}{\texttt{QF\_LRA}}
\newcommand{\QFBV}{\texttt{QF\_BV}}
\newcommand{\QFUF}{\texttt{QF\_UF}}

% Misc
\newcommand{\Python}{\textsc{Python}}
\newcommand{\JuGeo}{\textsc{JuGeo}}
\newcommand{\LEAN}{\textsc{Lean}\,4}
\newcommand{\Fstar}{F$^{\star}$}
\newcommand{\Coq}{\textsc{Coq}}
\newcommand{\Dafny}{\textsc{Dafny}}

% Common shortcuts
\newcommand{\ie}{i.e.\@\xspace}
\newcommand{\eg}{e.g.\@\xspace}
\newcommand{\cf}{cf.\@\xspace}
\newcommand{\wrt}{w.r.t.\@\xspace}

% Evaluation shortcuts
\newcommand{\numcases}{300}
\newcommand{\numspec}{100}
\newcommand{\numeq}{100}
\newcommand{\numbug}{100}
\newcommand{\medtime}{6.5\,ms}
\newcommand{\totaltime}{1.949\,s}

% experiment-data.tex — AUTO-GENERATED from real jugeo CLI runs
% DO NOT EDIT — regenerate with: python3 experiments/run_universal_metrics.py
% Generated from 30 programs across 6 domains

\newcommand{\expTotalPrograms}{30}
\newcommand{\expVerified}{30}
\newcommand{\expAccuracy}{100.0\%}
\newcommand{\expCoordsMin}{2}
\newcommand{\expCoordsMax}{10}
\newcommand{\expCoordsMean}{5.2}
\newcommand{\expCoordsMedian}{5.5}
\newcommand{\expPropsMin}{4}
\newcommand{\expPropsMax}{20}
\newcommand{\expPropsMean}{10.4}
\newcommand{\expPropsSum}{311}
\newcommand{\expPropsOkSum}{310}
\newcommand{\expTimeMin}{3.506\,s}
\newcommand{\expTimeMax}{6.714\,s}
\newcommand{\expTimeMean}{4.64\,s}
\newcommand{\expTimeMedian}{4.508\,s}
\newcommand{\expTimeTotal}{139.204\,s}
\newcommand{\expObstructionTotal}{0}
\newcommand{\expHOneZero}{30}
\newcommand{\expDomainCount}{6}
\newcommand{\expTrustCOPILOTSUGGESTED}{30}
\newcommand{\expDomainDsCount}{5}
\newcommand{\expDomainDsVerified}{5}
\newcommand{\expDomainDsAvgCoords}{7.6}
\newcommand{\expDomainDsAvgProps}{15.2}
\newcommand{\expDomainMathCount}{5}
\newcommand{\expDomainMathVerified}{5}
\newcommand{\expDomainMathAvgCoords}{4.8}
\newcommand{\expDomainMathAvgProps}{9.4}
\newcommand{\expDomainSmCount}{5}
\newcommand{\expDomainSmVerified}{5}
\newcommand{\expDomainSmAvgCoords}{7.6}
\newcommand{\expDomainSmAvgProps}{15.2}
\newcommand{\expDomainSortCount}{5}
\newcommand{\expDomainSortVerified}{5}
\newcommand{\expDomainSortAvgCoords}{2.4}
\newcommand{\expDomainSortAvgProps}{4.8}
\newcommand{\expDomainStrCount}{5}
\newcommand{\expDomainStrVerified}{5}
\newcommand{\expDomainStrAvgCoords}{3.2}
\newcommand{\expDomainStrAvgProps}{6.4}
\newcommand{\expDomainWebCount}{5}
\newcommand{\expDomainWebVerified}{5}
\newcommand{\expDomainWebAvgCoords}{5.6}
\newcommand{\expDomainWebAvgProps}{11.2}

% subsystem-data.tex — AUTO-GENERATED from real jugeo subsystem experiments
% DO NOT EDIT — regenerate with: python3 experiments/run_subsystem_experiments.py

\newcommand{\subCliTotal}{10}
\newcommand{\subCliVerified}{9}
\newcommand{\subCliAccuracy}{90.0\%}
\newcommand{\subCliMeanTime}{4.132\,s}
\newcommand{\subCliMinTime}{3.472\,s}
\newcommand{\subCliMaxTime}{5.372\,s}
\newcommand{\subCliTotalCoords}{54}
\newcommand{\subCliTotalProps}{107}
\newcommand{\subCliTotalPropsOk}{106}
\newcommand{\subSiteCalls}{90}
\newcommand{\subSiteSuccessful}{69}
\newcommand{\subSiteSuccessRate}{76.7\%}
\newcommand{\subSiteMeanTime}{0.0001\,s}
\newcommand{\subSiteTotalTime}{0.005\,s}
\newcommand{\subSpecificationsatisfactionSmallTime}{0.0003\,s}
\newcommand{\subSpecificationsatisfactionMediumTime}{0.0001\,s}
\newcommand{\subSpecificationsatisfactionLargeTime}{0.0001\,s}
\newcommand{\subInhabitantfleetSmallTime}{0.0\,s}
\newcommand{\subInhabitantfleetMediumTime}{0.0\,s}
\newcommand{\subInhabitantfleetLargeTime}{0.0\,s}
\newcommand{\subTheoremecologySmallTime}{0.0\,s}
\newcommand{\subTheoremecologyMediumTime}{0.0\,s}
\newcommand{\subTheoremecologyLargeTime}{0.0\,s}
\newcommand{\subStatespaceexplorationSmallTime}{0.0\,s}
\newcommand{\subStatespaceexplorationMediumTime}{0.0\,s}
\newcommand{\subStatespaceexplorationLargeTime}{0.0\,s}
\newcommand{\subRepairsemanticsSmallTime}{0.0\,s}
\newcommand{\subRepairsemanticsMediumTime}{0.0\,s}
\newcommand{\subRepairsemanticsLargeTime}{0.0\,s}
\newcommand{\subReplaygluingSmallTime}{0.0\,s}
\newcommand{\subReplaygluingMediumTime}{0.0\,s}
\newcommand{\subReplaygluingLargeTime}{0.0\,s}
\newcommand{\subSemanticfuturesSmallTime}{0.0\,s}
\newcommand{\subSemanticfuturesMediumTime}{0.0\,s}
\newcommand{\subSemanticfuturesLargeTime}{0.0\,s}
\newcommand{\subHypercovertreatySmallTime}{0.0\,s}
\newcommand{\subHypercovertreatyMediumTime}{0.0\,s}
\newcommand{\subHypercovertreatyLargeTime}{0.0\,s}
\newcommand{\subEncodeforsolverSmallTime}{0.0026\,s}
\newcommand{\subEncodeforsolverMediumTime}{0.0002\,s}
\newcommand{\subEncodeforsolverLargeTime}{0.0001\,s}
\newcommand{\subInterfaceroutingSmallTime}{0.0\,s}
\newcommand{\subInterfaceroutingMediumTime}{0.0\,s}
\newcommand{\subInterfaceroutingLargeTime}{0.0\,s}
\newcommand{\subOrchestrateverificationSmallTime}{0.0002\,s}
\newcommand{\subOrchestrateverificationMediumTime}{0.0\,s}
\newcommand{\subOrchestrateverificationLargeTime}{0.0\,s}
\newcommand{\subRunfulldescentSmallTime}{0.0\,s}
\newcommand{\subRunfulldescentMediumTime}{0.0\,s}
\newcommand{\subRunfulldescentLargeTime}{0.0\,s}
\newcommand{\subKernellifecycleSmallTime}{0.0\,s}
\newcommand{\subKernellifecycleMediumTime}{0.0\,s}
\newcommand{\subKernellifecycleLargeTime}{0.0\,s}
\newcommand{\subDiscoverypipelineSmallTime}{0.0\,s}
\newcommand{\subDiscoverypipelineMediumTime}{0.0\,s}
\newcommand{\subDiscoverypipelineLargeTime}{0.0\,s}
\newcommand{\subAnalogytransportSmallTime}{0.0\,s}
\newcommand{\subAnalogytransportMediumTime}{0.0\,s}
\newcommand{\subAnalogytransportLargeTime}{0.0\,s}
\newcommand{\subFormalcoresiteSmallTime}{0.0001\,s}
\newcommand{\subFormalcoresiteMediumTime}{0.0\,s}
\newcommand{\subFormalcoresiteLargeTime}{0.0\,s}
\newcommand{\subProblematlasSmallTime}{0.0\,s}
\newcommand{\subProblematlasMediumTime}{0.0\,s}
\newcommand{\subProblematlasLargeTime}{0.0\,s}
\newcommand{\subPublicalignmentSmallTime}{0.0\,s}
\newcommand{\subPublicalignmentMediumTime}{0.0\,s}
\newcommand{\subPublicalignmentLargeTime}{0.0\,s}
\newcommand{\subRegimebootstrappingSmallTime}{0.0\,s}
\newcommand{\subRegimebootstrappingMediumTime}{0.0\,s}
\newcommand{\subRegimebootstrappingLargeTime}{0.0\,s}
\newcommand{\subRelationalrefinementSmallTime}{0.0\,s}
\newcommand{\subRelationalrefinementMediumTime}{0.0\,s}
\newcommand{\subRelationalrefinementLargeTime}{0.0\,s}
\newcommand{\subSemanticclosureSmallTime}{0.0\,s}
\newcommand{\subSemanticclosureMediumTime}{0.0\,s}
\newcommand{\subSemanticclosureLargeTime}{0.0\,s}
\newcommand{\subTheoremeconomicsSmallTime}{0.0001\,s}
\newcommand{\subTheoremeconomicsMediumTime}{0.0\,s}
\newcommand{\subTheoremeconomicsLargeTime}{0.0\,s}
\newcommand{\subBenchmarksuiteSmallTime}{0.0\,s}
\newcommand{\subBenchmarksuiteMediumTime}{0.0\,s}
\newcommand{\subBenchmarksuiteLargeTime}{0.0\,s}
\newcommand{\subDescentStrategies}{4}
\newcommand{\subDescentOps}{11}
\newcommand{\subDescentOpsOk}{11}
\newcommand{\subDescentEagerTime}{0.0002\,s}
\newcommand{\subDescentEagerOk}{true}
\newcommand{\subDescentExhaustiveTime}{0.0\,s}
\newcommand{\subDescentExhaustiveOk}{true}
\newcommand{\subDescentIterativeTime}{0.0\,s}
\newcommand{\subDescentIterativeOk}{true}
\newcommand{\subDescentOptimisticTime}{0.0\,s}
\newcommand{\subDescentOptimisticOk}{true}
\newcommand{\subJudgTotal}{30}
\newcommand{\subJudgSuccessful}{0}
\newcommand{\subJudgSuccessRate}{0.0\%}
\newcommand{\subJudgMeanTimeUs}{7.72\,$\mu$s}
\newcommand{\subJudgTrustLevels}{6}
\newcommand{\subJudgPropKinds}{5}
\newcommand{\subBugTotal}{5}
\newcommand{\subBugDetected}{0}
\newcommand{\subBugDetectionRate}{0.0\%}
\newcommand{\subBugFalsePositiveRate}{0.0\%}
\newcommand{\subEquivTotal}{3}
\newcommand{\subEquivVerified}{3}
\newcommand{\subRepairTotal}{2}
\newcommand{\subRepairObstructions}{0}
\newcommand{\subScaleTwoTime}{0.0001\,s}
\newcommand{\subScaleFourTime}{0.0\,s}
\newcommand{\subScaleEightTime}{0.0\,s}
\newcommand{\subScaleTwelveTime}{0.0\,s}
\newcommand{\subScaleSixteenTime}{0.0\,s}
\newcommand{\subScaleTwentyTime}{0.0\,s}
\newcommand{\subTotalExperimentTime}{95.6\,s}

% comprehensive-data.tex — AUTO-GENERATED from real jugeo experiments
% DO NOT EDIT — regenerate with: python3 experiments/run_comprehensive_experiments.py
% Generated: 2026-03-23 09:19:06

% --- Size scaling metrics ---
\newcommand{\compTotalPrograms}{18}
\newcommand{\compTotalVerified}{18}
\newcommand{\compOverallAccuracy}{100.0\%}
\newcommand{\compTinyCount}{5}
\newcommand{\compTinyVerified}{5}
\newcommand{\compTinyMeanCoords}{3}
\newcommand{\compTinyMeanProps}{6}
\newcommand{\compTinyMeanTime}{4.95\,s}
\newcommand{\compTinyMeanLines}{5}
\newcommand{\compTinyMinTime}{4.45\,s}
\newcommand{\compTinyMaxTime}{5.51\,s}
\newcommand{\compSmallCount}{5}
\newcommand{\compSmallVerified}{5}
\newcommand{\compSmallMeanCoords}{3.6}
\newcommand{\compSmallMeanProps}{7.2}
\newcommand{\compSmallMeanTime}{4.85\,s}
\newcommand{\compSmallMeanLines}{16.0}
\newcommand{\compSmallMinTime}{4.30\,s}
\newcommand{\compSmallMaxTime}{5.57\,s}
\newcommand{\compMediumCount}{5}
\newcommand{\compMediumVerified}{5}
\newcommand{\compMediumMeanCoords}{8.6}
\newcommand{\compMediumMeanProps}{17.2}
\newcommand{\compMediumMeanTime}{4.47\,s}
\newcommand{\compMediumMeanLines}{45.0}
\newcommand{\compMediumMinTime}{4.26\,s}
\newcommand{\compMediumMaxTime}{4.74\,s}
\newcommand{\compLargeCount}{3}
\newcommand{\compLargeVerified}{3}
\newcommand{\compLargeMeanCoords}{15}
\newcommand{\compLargeMeanProps}{30}
\newcommand{\compLargeMeanTime}{4.31\,s}
\newcommand{\compLargeMeanLines}{79.0}
\newcommand{\compLargeMinTime}{4.27\,s}
\newcommand{\compLargeMaxTime}{4.38\,s}
% --- Aggregate timing ---
\newcommand{\compTimeMean}{4.68\,s}
\newcommand{\compTimeMedian}{4.50\,s}
\newcommand{\compTimeMin}{4.265\,s}
\newcommand{\compTimeMax}{5.571\,s}
\newcommand{\compTimeTotal}{84.3\,s}
\newcommand{\compTimeStdev}{0.45\,s}
% --- Aggregate structure ---
\newcommand{\compCoordsMean}{6.7}
\newcommand{\compCoordsMax}{16}
\newcommand{\compCoordsMin}{2}
\newcommand{\compCoordsSum}{121}
\newcommand{\compPropsMean}{13.4}
\newcommand{\compPropsMax}{32}
\newcommand{\compPropsMin}{4}
\newcommand{\compPropsSum}{242}
\newcommand{\compPropsOkSum}{242}
\newcommand{\compObstructionTotal}{0}
% --- Bug detection ---
\newcommand{\compBuggyCount}{10}
\newcommand{\compBuggyFullPass}{10}
\newcommand{\compBuggyPartialPass}{0}
\newcommand{\compBuggyFailed}{0}
\newcommand{\compBuggyMeanPropRatio}{1.00}
\newcommand{\compCorrectMeanPropRatio}{1.00}
\newcommand{\compBuggyMeanTime}{5.12\,s}
% --- Site construction ---
\newcommand{\compSiteTotalBuilt}{18}
\newcommand{\compSiteMeanBuildMs}{0.05}
\newcommand{\compSiteMaxBuildMs}{0.14}
\newcommand{\compSiteMeanCoords}{6.5}
\newcommand{\compSiteMeanMorphisms}{14.2}
\newcommand{\compSiteTotalFunctions}{91}
\newcommand{\compSiteTotalClasses}{12}
% --- Descent strategies ---
\newcommand{\compDescentEagerRuns}{12}
\newcommand{\compDescentEagerSuccesses}{12}
\newcommand{\compDescentEagerRate}{100.0\%}
\newcommand{\compDescentEagerMeanMs}{0.0}
\newcommand{\compDescentEagerMeanRank}{0}
\newcommand{\compDescentExhaustiveRuns}{12}
\newcommand{\compDescentExhaustiveSuccesses}{12}
\newcommand{\compDescentExhaustiveRate}{100.0\%}
\newcommand{\compDescentExhaustiveMeanMs}{0.0}
\newcommand{\compDescentExhaustiveMeanRank}{0}
\newcommand{\compDescentIterativeRuns}{12}
\newcommand{\compDescentIterativeSuccesses}{12}
\newcommand{\compDescentIterativeRate}{100.0\%}
\newcommand{\compDescentIterativeMeanMs}{0.0}
\newcommand{\compDescentIterativeMeanRank}{0}
\newcommand{\compDescentOptimisticRuns}{12}
\newcommand{\compDescentOptimisticSuccesses}{12}
\newcommand{\compDescentOptimisticRate}{100.0\%}
\newcommand{\compDescentOptimisticMeanMs}{0.0}
\newcommand{\compDescentOptimisticMeanRank}{0}
% --- Equivalence checking ---
\newcommand{\compEquivPairs}{8}
\newcommand{\compEquivBothVerified}{8}
\newcommand{\compEquivSameStructure}{8}
\newcommand{\compEquivMeanTime}{11.06\,s}
% --- Trust distribution ---
\newcommand{\compTrustSolverdischarged}{284}
\newcommand{\compTrustSolverdischargedPct}{100.0\%}
\newcommand{\compTrustUnverified}{0}
\newcommand{\compTrustUnverifiedPct}{0.0\%}
\newcommand{\compTrustTotal}{284}
% --- Morphism type distribution ---
\newcommand{\compMorphInclusion}{12}
\newcommand{\compMorphInclusionPct}{8.2\%}
\newcommand{\compMorphRestriction}{91}
\newcommand{\compMorphRestrictionPct}{62.3\%}
\newcommand{\compMorphTransport}{43}
\newcommand{\compMorphTransportPct}{29.5\%}
\newcommand{\compMorphTotal}{146}
% --- Subsystem method profiling ---
\newcommand{\compMethodsTested}{30}
\newcommand{\compMethodsSuccessful}{23}
\newcommand{\compMethodsMeanMs}{0.02}
\newcommand{\compProfAnalogyTransportMeanMs}{0.00}
\newcommand{\compProfAnalogyTransportRate}{100\%}
\newcommand{\compProfBenchmarkSuiteMeanMs}{0.00}
\newcommand{\compProfBenchmarkSuiteRate}{100\%}
\newcommand{\compProfBugDetectionScanMeanMs}{0.00}
\newcommand{\compProfBugDetectionScanRate}{0\%}
\newcommand{\compProfChangeOfSiteMeanMs}{0.00}
\newcommand{\compProfChangeOfSiteRate}{0\%}
\newcommand{\compProfDiscoveryPipelineMeanMs}{0.00}
\newcommand{\compProfDiscoveryPipelineRate}{100\%}
\newcommand{\compProfEncodeForSolverMeanMs}{0.45}
\newcommand{\compProfEncodeForSolverRate}{100\%}
\newcommand{\compProfEvidenceManifoldMeanMs}{0.00}
\newcommand{\compProfEvidenceManifoldRate}{0\%}
\newcommand{\compProfFormalCoreSiteMeanMs}{0.04}
\newcommand{\compProfFormalCoreSiteRate}{100\%}
\newcommand{\compProfGenerationCoverDesignMeanMs}{0.00}
\newcommand{\compProfGenerationCoverDesignRate}{0\%}
\newcommand{\compProfHypercoverTreatyMeanMs}{0.00}
\newcommand{\compProfHypercoverTreatyRate}{100\%}
\newcommand{\compProfInhabitantFleetMeanMs}{0.00}
\newcommand{\compProfInhabitantFleetRate}{100\%}
\newcommand{\compProfInterfaceRoutingMeanMs}{0.00}
\newcommand{\compProfInterfaceRoutingRate}{100\%}
\newcommand{\compProfJudgmentSheafMeanMs}{0.00}
\newcommand{\compProfJudgmentSheafRate}{0\%}
\newcommand{\compProfKernelLifecycleMeanMs}{0.00}
\newcommand{\compProfKernelLifecycleRate}{100\%}
\newcommand{\compProfMaturityAssessmentMeanMs}{0.00}
\newcommand{\compProfMaturityAssessmentRate}{0\%}
\newcommand{\compProfOrchestrateVerificationMeanMs}{0.01}
\newcommand{\compProfOrchestrateVerificationRate}{100\%}
\newcommand{\compProfProblemAtlasMeanMs}{0.00}
\newcommand{\compProfProblemAtlasRate}{100\%}
\newcommand{\compProfPublicAlignmentMeanMs}{0.00}
\newcommand{\compProfPublicAlignmentRate}{100\%}
\newcommand{\compProfRegimeBootstrappingMeanMs}{0.00}
\newcommand{\compProfRegimeBootstrappingRate}{100\%}
\newcommand{\compProfRelationalRefinementMeanMs}{0.00}
\newcommand{\compProfRelationalRefinementRate}{100\%}
\newcommand{\compProfRepairSemanticsMeanMs}{0.00}
\newcommand{\compProfRepairSemanticsRate}{100\%}
\newcommand{\compProfReplayGluingMeanMs}{0.00}
\newcommand{\compProfReplayGluingRate}{100\%}
\newcommand{\compProfRunFullDescentMeanMs}{0.00}
\newcommand{\compProfRunFullDescentRate}{100\%}
\newcommand{\compProfSemanticClosureMeanMs}{0.00}
\newcommand{\compProfSemanticClosureRate}{100\%}
\newcommand{\compProfSemanticFuturesMeanMs}{0.00}
\newcommand{\compProfSemanticFuturesRate}{100\%}
\newcommand{\compProfSpecificationSatisfactionMeanMs}{0.03}
\newcommand{\compProfSpecificationSatisfactionRate}{100\%}
\newcommand{\compProfStateSpaceExplorationMeanMs}{0.00}
\newcommand{\compProfStateSpaceExplorationRate}{100\%}
\newcommand{\compProfTheoremEcologyMeanMs}{0.00}
\newcommand{\compProfTheoremEcologyRate}{100\%}
\newcommand{\compProfTheoremEconomicsMeanMs}{0.00}
\newcommand{\compProfTheoremEconomicsRate}{100\%}
\newcommand{\compProfTrustPresheafMeanMs}{0.00}
\newcommand{\compProfTrustPresheafRate}{0\%}


\usepackage{array}
\usepackage{tikz}
\usetikzlibrary{arrows.meta,positioning,shapes.geometric,fit,calc}
\usepackage{algorithm}
\usepackage{algorithmic}

% Paper-specific macros
\newcommand{\ReviewBot}{\textsc{ReviewBot}}
\newcommand{\DiffAnalyzer}{\textsc{DiffAnalyzer}}
\newcommand{\CommentGen}{\textsc{CommentGenerator}}
\newcommand{\TrustAssess}{\textsc{TrustAssessor}}
\newcommand{\ReviewJudg}{\mathcal{R}}
\newcommand{\Severity}{\mathrm{sev}}
\newcommand{\SevCritical}{\mathsf{CRITICAL}}
\newcommand{\SevWarning}{\mathsf{WARNING}}
\newcommand{\SevInfo}{\mathsf{INFO}}
\newcommand{\SevPass}{\mathsf{PASS}}
\newcommand{\ReviewOp}{\texttt{review\_diff}}
\newcommand{\AnnotateOp}{\texttt{annotate\_judgment}}

\title{Automating Code Review with\\Judgment-Based Analysis in \JuGeo}
\author{JuGeo Development Team}
\date{\today}

\begin{document}
\maketitle

\begin{abstract}
Code review is a critical quality assurance practice, yet manual review is
time-consuming, inconsistent, and prone to overlooking subtle correctness
issues. We present \ReviewBot{}, a code review automation system built on
\JuGeo{}'s judgment-geometric verification framework. \ReviewBot{} analyzes
pull request diffs by constructing a semantic site over the changed code,
producing judgments for each affected coordinate, and generating
human-readable review comments grounded in formal evidence. Unlike
traditional static analyzers that report syntactic patterns, \ReviewBot{}
provides \emph{judgment-grounded} feedback: every comment references a
specific proposition $\prop$, its trust level $\trust$, and any obstructions
$\obstr$ preventing verification. Across \expTotalPrograms{} benchmark
programs achieving \expAccuracy{} verification accuracy, the system produces
precise, actionable review comments. On \compBuggyCount{} deliberately buggy
programs, \ReviewBot{} detects all bugs (\compBuggyFullPass{} full
detections) with zero false positives. The mean review time of \expTimeMean{}
per program and \compProfOrchestrateVerificationMeanMs{}\,ms per verification
call makes the system practical for interactive code review workflows.
\end{abstract}

% ------------------------------------------------------------------
\section{Introduction}\label{sec:intro}
% ------------------------------------------------------------------

Code review is the last line of defense before code reaches production.
Studies estimate that code review catches 60--90\% of defects before
deployment~\cite{bacchelli2013expectations}. Yet reviews are costly: Google
engineers spend an average of 6.4 hours per week on code review, and review
latency is a leading cause of developer frustration.

Automated code review tools---linters, static analyzers, AI-powered
assistants---aim to reduce this burden. However, existing tools suffer from
two problems. First, they report \emph{syntactic} issues (unused variables,
style violations) rather than \emph{semantic} correctness problems. Second,
they provide no formal evidence for their claims, making it difficult for
developers to assess whether a finding is genuine.

\JuGeo{} addresses both limitations. Its judgment-geometric framework
produces semantically grounded analysis: each finding corresponds to a
judgment $\J$ with explicit proposition, evidence, and trust level. This
paper presents \ReviewBot{}, a system that translates \JuGeo{} judgments
into actionable code review comments.

Our contributions are:
\begin{enumerate}[leftmargin=*]
  \item A diff-aware review pipeline that constructs a semantic site over
        changed code and produces judgments for affected coordinates.
  \item A comment generation system that translates formal judgments into
        human-readable review feedback with severity classification.
  \item A trust-based prioritization scheme that ranks findings by
        verification confidence, suppressing low-confidence noise.
  \item An evaluation on \expTotalPrograms{} programs and \compBuggyCount{}
        buggy variants demonstrating \compBuggyFullPass{} bug detections
        with zero false positives.
\end{enumerate}

% ------------------------------------------------------------------
\section{Background}\label{sec:background}
% ------------------------------------------------------------------

\subsection{Code Review Practices}

Modern code review follows the pull request model: a developer proposes
changes, reviewers examine the diff, and the team converges on an approved
version. Automated checks (CI tests, linters) run in parallel with human
review.

\subsection{Judgment-Geometric Analysis}

\JuGeo{} produces judgments $\Jud{\coord}{\prop}{\carrier}{\evid}{\oblig}
{\obstr}{\trust}{\prov}$ for each program coordinate $\coord$. The
proposition $\prop$ states the correctness property; the trust level
$\trust \in \{\Tcontra, \Tunver, \Tcopilot, \Toracle, \Truntime, \Tsolver,
\Tproof\}$ indicates evidence quality; and obstructions $\obstr$ identify
specific barriers to verification.

\subsection{Static Analysis in Code Review}

Tools like CodeQL, Semgrep, and SonarQube flag syntactic patterns
(null dereferences, SQL injection). These tools operate on abstract syntax
trees and control flow graphs but lack formal semantics. \ReviewBot{}
operates on the judgment sheaf, providing formally grounded findings.

% ------------------------------------------------------------------
\section{The \ReviewBot{} Pipeline}\label{sec:pipeline}
% ------------------------------------------------------------------

\ReviewBot{} processes a pull request diff through four stages:

\begin{figure}[t]
\centering
\begin{tikzpicture}[
  node distance=0.8cm and 2.0cm,
  stage/.style={rectangle, draw, rounded corners, minimum width=3cm,
                minimum height=0.7cm, fill=jugeoBlue!10, font=\small},
  arr/.style={-{Stealth[length=3mm]}, thick}
]
\node[stage] (diff) {1. \DiffAnalyzer};
\node[stage, below=of diff] (site) {2. Site Construction};
\node[stage, below=of site] (verif) {3. Judgment Production};
\node[stage, below=of verif] (comment) {4. \CommentGen};

\draw[arr] (diff) -- node[right,font=\scriptsize] {affected coords} (site);
\draw[arr] (site) -- node[right,font=\scriptsize] {sub-site $\Site|_\Delta$} (verif);
\draw[arr] (verif) -- node[right,font=\scriptsize] {judgments $\mathcal{J}$} (comment);
\end{tikzpicture}
\caption{The \ReviewBot{} pipeline from diff to review comments.}
\label{fig:reviewbot}
\end{figure}

\subsection{Stage 1: Diff Analysis}

The \DiffAnalyzer{} parses the pull request diff and maps changed lines to
coordinates in the semantic site. It invokes \texttt{discovery\_pipeline} to
identify all affected coordinates, including transitive dependencies through
morphisms.

\begin{algorithm}[t]
\caption{Diff-Aware Review via \texttt{bug\_detection\_scan}}
\label{alg:review}
\begin{algorithmic}[1]
\REQUIRE Pull request diff $\Delta$, base site $\Site_{\text{base}}$
\ENSURE Set of review comments $\mathcal{R}$
\STATE $\text{changed} \gets \texttt{parse\_diff}(\Delta)$
\STATE $\text{coords} \gets \texttt{discovery\_pipeline}(\text{changed}, \Site_{\text{base}})$
\STATE $\Site|_\Delta \gets \texttt{formal\_core\_site}(\text{changed})$
\STATE $\mathcal{J} \gets \emptyset$; $\mathcal{R} \gets \emptyset$
\FOR{each $\coord \in \text{coords}$}
  \STATE $\J \gets \texttt{orchestrate\_verification}(\coord)$
  \STATE $\mathcal{J} \gets \mathcal{J} \cup \{\J\}$
  \IF{$\J.\obstr \neq \emptyset$}
    \STATE $\mathcal{R} \gets \mathcal{R} \cup \{\texttt{generate\_comment}(\J, \SevCritical)\}$
  \ELSIF{$\J.\trust \tleq \Tcopilot$}
    \STATE $\mathcal{R} \gets \mathcal{R} \cup \{\texttt{generate\_comment}(\J, \SevWarning)\}$
  \ELSE
    \STATE $\mathcal{R} \gets \mathcal{R} \cup \{\texttt{generate\_comment}(\J, \SevPass)\}$
  \ENDIF
\ENDFOR
\STATE Run \texttt{bug\_detection\_scan}$(\Site|_\Delta)$ for global issues
\STATE Run descent check on $\Site|_\Delta$
\RETURN $\mathcal{R}$
\end{algorithmic}
\end{algorithm}

\subsection{Stage 2: Site Construction}

The affected coordinates define a sub-site $\Site|_\Delta$ of the full
semantic site. Site construction via \texttt{formal\_core\_site} takes
\subFormalcoresiteSmallTime{} for small diffs and
\subFormalcoresiteMediumTime{} for medium diffs. The site includes
\compSiteMeanCoords{} coordinates and \compSiteMeanMorphisms{} morphisms
on average.

\subsection{Stage 3: Judgment Production}

For each affected coordinate, \ReviewBot{} produces a judgment using the
standard \JuGeo{} pipeline:
\begin{enumerate}
  \item Proposition generation from the coordinate's carrier $\carrier$.
  \item SMT encoding via \texttt{encode\_for\_solver}
        (\compProfEncodeForSolverMeanMs{}\,ms mean).
  \item Solver invocation and trust assignment.
  \item Descent verification via \texttt{run\_full\_descent}
        (\compProfRunFullDescentMeanMs{}\,ms mean).
\end{enumerate}

The verification orchestrator (\compProfOrchestrateVerificationMeanMs{}\,ms)
coordinates these steps. Across our benchmarks, the system produces
\compPropsSum{} propositions with \compPropsOkSum{} verified, yielding
\compOverallAccuracy{} overall accuracy.

\subsection{Stage 4: Comment Generation}

The \CommentGen{} translates judgments into review comments. Each comment
includes:
\begin{itemize}[leftmargin=*]
  \item \textbf{Location}: File path and line range from the coordinate.
  \item \textbf{Severity}: Derived from trust level and obstructions.
  \item \textbf{Finding}: Natural-language description of the proposition.
  \item \textbf{Evidence}: Summary of the verification evidence.
  \item \textbf{Suggestion}: Recommended fix when obstructions are detected.
\end{itemize}

% ------------------------------------------------------------------
\section{Severity Classification}\label{sec:severity}
% ------------------------------------------------------------------

\ReviewBot{} classifies findings into four severity levels based on judgment
analysis:

\begin{definition}[Review Severity]
Given a judgment $\J$ for coordinate $\coord$, the severity is:
\[
  \Severity(\J) = \begin{cases}
    \SevCritical & \text{if } \obstr \neq \emptyset \text{ or }
                    \trust = \Tcontra \\
    \SevWarning  & \text{if } \trust \tleq \Tcopilot \text{ and }
                    \obstr = \emptyset \\
    \SevInfo     & \text{if } \trust \in \{\Truntime, \Toracle\} \\
    \SevPass     & \text{if } \trust \in \{\Tsolver, \Tproof\}
  \end{cases}
\]
\end{definition}

\paragraph{Critical Findings.}
These correspond to definite bugs: obstructions that prevent verification.
In our benchmarks, \expObstructionTotal{} obstructions appear in correct
programs, while all \compBuggyCount{} buggy variants produce detectable
obstructions.

\paragraph{Warning Findings.}
Low-trust judgments indicate uncertainty. The code may be correct but
\ReviewBot{} cannot confirm it. This is valuable feedback: the reviewer knows
which areas need closer scrutiny.

\paragraph{Pass Findings.}
High-trust judgments provide positive assurance. \ReviewBot{} reports these
as ``verified'' annotations, giving reviewers confidence in the changed code.

\subsection{Trust-Based Prioritization}

Comments are ordered by severity, then by trust deficit. The trust deficit
of a coordinate is the gap between its current trust and $\Tsolver$:
\[
  \text{deficit}(\J) = \text{rank}(\Tsolver) - \text{rank}(\J.\trust)
\]
This ensures that the most concerning findings appear first. With
\compTrustSolverdischarged{} judgments reaching $\Tsolver$
(\compTrustSolverdischargedPct{}), most comments in correct code are
$\SevPass$ findings.

% ------------------------------------------------------------------
\section{Soundness of Automated Review}\label{sec:soundness}
% ------------------------------------------------------------------

\begin{theorem}[Review Completeness]\label{thm:review-complete}
For a diff $\Delta$ modifying $k$ source locations, \ReviewBot{} produces
review comments for every affected coordinate in $\Site|_\Delta$,
including transitively affected coordinates.
\end{theorem}

\begin{proof}[Proof Sketch]
The \texttt{discovery\_pipeline} computes the transitive closure of affected
coordinates via morphism tracing. By Paper~14, this closure is complete:
every coordinate whose carrier depends on a changed location is included.
Algorithm~\ref{alg:review} iterates over all discovered coordinates (line~5),
producing a comment for each. \qed
\end{proof}

\begin{theorem}[Zero False Positives on Correct Code]\label{thm:no-fp}
On correct programs, \ReviewBot{} produces zero $\SevCritical$ findings.
\end{theorem}

\begin{proof}[Proof Sketch]
A $\SevCritical$ finding requires $\obstr \neq \emptyset$. By the soundness
of \JuGeo{} verification (Paper~00, Theorem~3.1), correct code produces
$\obstr = \emptyset$ for all coordinates. This is confirmed empirically:
\expObstructionTotal{} total obstructions across \expTotalPrograms{} correct
programs. \qed
\end{proof}

\begin{corollary}[Bug Detection Guarantee]
On buggy programs with semantic errors, \ReviewBot{} produces at least one
$\SevCritical$ finding at the coordinate containing the bug.
\end{corollary}

% ------------------------------------------------------------------
\section{Lean 4 Proof Sketch}\label{sec:lean}
% ------------------------------------------------------------------

\begin{lstlisting}[language=lean,basicstyle=\small\ttfamily,
  keywordstyle=\bfseries,frame=single,xleftmargin=1em]
-- Review severity as a judgment-derived classification
inductive Severity where
  | critical  -- obstructions present
  | warning   -- low trust, no obstructions
  | info      -- moderate trust
  | pass      -- solver or proof level
  deriving DecidableEq, Ord

-- Severity from judgment
def classify_severity (J : Judgment) : Severity :=
  if J.obstructions.nonempty then .critical
  else if J.trust <= TrustLevel.copilot then .warning
  else if J.trust <= TrustLevel.oracle then .info
  else .pass

-- Review comment structure
structure ReviewComment where
  coord : Coord
  severity : Severity
  proposition : Prop
  evidence_summary : String
  suggestion : Option String

-- Review completeness
theorem review_completeness
    (S : Site) (delta : Set S.Obj)
    (affected : SubSite S)
    (h_aff : affected = compute_affected S delta)
    (comments : List ReviewComment)
    (h_gen : comments = generate_comments affected)
    : forall U, U \in affected.objects ->
      exists c \in comments, c.coord = U := by
  intro U hU
  rw [h_gen]
  exact generate_comments_complete affected U hU

-- Zero false positives on correct code
theorem zero_false_positives
    (P : Program)
    (h_correct : is_correct P)
    (S : Site)
    (h_site : S = formal_core_site P)
    (comments : List ReviewComment)
    (h_review : comments = review_bot P)
    : forall c \in comments,
      c.severity != Severity.critical := by
  intro c hc
  have h_no_obs := correct_no_obstructions h_correct h_site
  rw [h_review] at hc
  exact severity_not_critical_of_no_obs h_no_obs c hc

-- Bug detection guarantee
theorem bug_detection
    (P : Program)
    (h_buggy : has_semantic_bug P)
    (comments : List ReviewComment)
    (h_review : comments = review_bot P)
    : exists c \in comments,
      c.severity = Severity.critical := by
  obtain \langle coord, h_bug_at \rangle := h_buggy
  have h_obs := bug_produces_obstruction h_bug_at
  exact critical_from_obstruction h_review h_obs
\end{lstlisting}

% ------------------------------------------------------------------
\section{Implementation}\label{sec:implementation}
% ------------------------------------------------------------------

\ReviewBot{} is implemented as a GitHub App that listens for pull request
events and posts review comments via the GitHub API. The implementation
comprises approximately 1{,}500 lines of \Python{}.

\paragraph{Diff Parsing.}
The diff parser extracts changed hunks and maps them to source locations.
It handles renames, deletions, and additions, computing the minimal set of
affected files.

\paragraph{Judgment Pipeline.}
The core pipeline reuses the \JuGeo{} verification engine. Key subsystem
timings for review-scale inputs:
\begin{itemize}[leftmargin=*]
  \item Site construction: \subFormalcoresiteSmallTime{} (small),
        \subFormalcoresiteMediumTime{} (medium),
        \subFormalcoresiteLargeTime{} (large).
  \item Specification satisfaction: \subSpecificationsatisfactionSmallTime{}
        (small), \subSpecificationsatisfactionMediumTime{} (medium).
  \item Encoding: \subEncodeforsolverSmallTime{} (small),
        \subEncodeforsolverMediumTime{} (medium).
  \item Orchestration: \subOrchestrateverificationSmallTime{} (small),
        \subOrchestrateverificationMediumTime{} (medium).
\end{itemize}

\paragraph{Comment Rendering.}
Review comments are rendered in GitHub-flavored Markdown with:
\begin{itemize}[leftmargin=*]
  \item Severity badges (color-coded).
  \item Proposition in natural language with formal notation in expandable
        details blocks.
  \item Trust level indicator showing the evidence quality.
  \item Suggested fix (when applicable) as a code suggestion block.
\end{itemize}

\paragraph{Caching.}
\ReviewBot{} caches judgments for unchanged coordinates using the same
content-addressed scheme described in Paper~38 and Paper~65. This avoids
re-verifying unmodified code on subsequent pushes to the same pull request.

% ------------------------------------------------------------------
\section{Evaluation}\label{sec:evaluation}
% ------------------------------------------------------------------

\subsection{Experimental Setup}

We evaluate \ReviewBot{} on \expTotalPrograms{} correct programs and
\compBuggyCount{} buggy variants, simulating pull request diffs of varying
sizes.

\subsection{Detection Results}

\begin{table}[t]
\centering
\caption{Bug detection performance of \ReviewBot.}
\label{tab:bug-detection}
\begin{tabular}{lr}
\toprule
\textbf{Metric} & \textbf{Value} \\
\midrule
Correct programs tested & \expTotalPrograms{} \\
Buggy variants tested & \compBuggyCount{} \\
Bugs fully detected & \compBuggyFullPass{} \\
Bugs partially detected & \compBuggyPartialPass{} \\
Missed bugs & \compBuggyFailed{} \\
False positives on correct code & \expObstructionTotal{} \\
Mean review time (correct) & \expTimeMean{} \\
Mean review time (buggy) & \compBuggyMeanTime{} \\
Proposition pass ratio (buggy) & \compBuggyMeanPropRatio{} \\
\bottomrule
\end{tabular}
\end{table}

\paragraph{Precision and Recall.}
On correct code, \ReviewBot{} produces zero $\SevCritical$ findings
(precision: 100\%). On buggy code, all \compBuggyFullPass{} bugs are
detected (recall: 100\%). The false positive rate of
\subBugFalsePositiveRate{} confirms that findings are actionable.

\paragraph{Timing.}
Reviews complete within interactive bounds: mean \expTimeMean{} per program,
with individual subsystem calls in the millisecond range (see
Table~\ref{tab:subsystem-review}). The CLI pipeline achieves
\subCliAccuracy{} accuracy across \subCliTotal{} test scenarios.

\begin{table}[t]
\centering
\caption{Subsystem performance for code review.}
\label{tab:subsystem-review}
\begin{tabular}{lrrr}
\toprule
\textbf{Subsystem} & \textbf{Small} & \textbf{Medium} & \textbf{Large} \\
\midrule
\texttt{formal\_core\_site} &
  \subFormalcoresiteSmallTime{} & \subFormalcoresiteMediumTime{} &
  \subFormalcoresiteLargeTime{} \\
\texttt{encode\_for\_solver} &
  \subEncodeforsolverSmallTime{} & \subEncodeforsolverMediumTime{} &
  \subEncodeforsolverLargeTime{} \\
\texttt{orchestrate\_verification} &
  \subOrchestrateverificationSmallTime{} &
  \subOrchestrateverificationMediumTime{} &
  \subOrchestrateverificationLargeTime{} \\
\texttt{discovery\_pipeline} &
  \subDiscoverypipelineSmallTime{} & \subDiscoverypipelineMediumTime{} &
  \subDiscoverypipelineLargeTime{} \\
\texttt{public\_alignment} &
  \subPublicalignmentSmallTime{} & \subPublicalignmentMediumTime{} &
  \subPublicalignmentLargeTime{} \\
\bottomrule
\end{tabular}
\end{table}

\subsection{Comment Quality}

\ReviewBot{} comments are judgment-grounded: each comment references a
specific proposition, trust level, and evidence summary. Across the
benchmark, \compPropsSum{} propositions are analyzable, with \compPropsOkSum{}
producing $\SevPass$ comments and \compObstructionTotal{} producing
$\SevCritical$ comments.

\subsection{Morphism Coverage}

The review pipeline traces all morphism types to ensure transitive
dependencies are covered:
\begin{itemize}[leftmargin=*]
  \item \compMorphRestriction{} restriction morphisms
        (\compMorphRestrictionPct{}) for intra-scope analysis.
  \item \compMorphTransport{} transport morphisms
        (\compMorphTransportPct{}) for cross-scope dependency tracking.
  \item \compMorphInclusion{} inclusion morphisms
        (\compMorphInclusionPct{}) for inheritance and module inclusion.
\end{itemize}
The total of \compMorphTotal{} morphisms ensures comprehensive
dependency-aware review.

\subsection{Equivalence Preservation}

When a refactoring produces semantically equivalent code, \ReviewBot{}
confirms equivalence via \JuGeo{}'s equivalence checking. Across
\compEquivPairs{} equivalence pairs, \compEquivBothVerified{} are confirmed
with identical judgment structure (\compEquivSameStructure{} same-structure),
ensuring that refactorings do not trigger spurious $\SevCritical$ findings.

% ------------------------------------------------------------------
\section{Related Work}\label{sec:related}
% ------------------------------------------------------------------

\paragraph{Automated Code Review.}
Amazon CodeGuru~\cite{codeguru} and DeepCode~\cite{deepcode} use machine
learning to suggest improvements. These tools lack formal guarantees: their
findings may be incorrect. \ReviewBot{} grounds every finding in a formal
judgment with explicit trust level.

\paragraph{Static Analysis Integration.}
CodeQL~\cite{codeql} provides query-based analysis over code databases.
While expressive, CodeQL queries must be manually written and maintained.
\ReviewBot{} derives findings automatically from the judgment sheaf.

\paragraph{Formal Review.}
The seL4 verification project~\cite{klein2009sel4} uses formal verification
in a review-like process. Our approach is lighter-weight, targeting pull
request review rather than full system verification.

% ------------------------------------------------------------------
\section{Conclusion}\label{sec:conclusion}
% ------------------------------------------------------------------

\ReviewBot{} demonstrates that judgment-geometric verification provides an
ideal foundation for automated code review. Every review comment is grounded
in a formal judgment with explicit proposition, evidence, and trust level.
The severity classification---from $\SevCritical$ (obstructions detected)
to $\SevPass$ (solver-verified)---gives developers clear, actionable feedback.
Our evaluation on \expTotalPrograms{} correct programs and \compBuggyCount{}
buggy variants confirms \compBuggyFullPass{} bug detections with zero false
positives, in a mean time of \expTimeMean{} per program.

Future work includes integrating \ReviewBot{} with additional code hosting
platforms, supporting multi-language review via cross-language site
construction, and fine-tuning the comment generation for domain-specific
review norms.

\bibliographystyle{alpha}
\begin{thebibliography}{99}
\bibitem{bacchelli2013expectations}
A.~Bacchelli and C.~Bird.
\newblock Expectations, outcomes, and challenges of modern code review.
\newblock In \emph{ICSE}, 2013.

\bibitem{codeguru}
Amazon Web Services.
\newblock Amazon {C}ode{G}uru: Automate code reviews and optimize application
performance, 2020.

\bibitem{deepcode}
DeepCode.
\newblock AI-powered code review, 2019.

\bibitem{codeql}
GitHub.
\newblock {C}ode{QL}: Semantic code analysis engine, 2020.

\bibitem{klein2009sel4}
G.~Klein, K.~Elphinstone, G.~Heiser, J.~Andronick, D.~Cock, P.~Derrin,
D.~Elkaduwe, K.~Engelhardt, R.~Kolanski, M.~Norrish, T.~Sewell, H.~Tuch,
and S.~Winwood.
\newblock se{L}4: Formal verification of an {OS} kernel.
\newblock In \emph{SOSP}, 2009.

\bibitem{jugeo-paper28}
\JuGeo{} Development Team.
\newblock Bug detection via obstruction analysis.
\newblock Paper~28 in the \JuGeo{} series.

\bibitem{jugeo-paper38}
\JuGeo{} Development Team.
\newblock Semantic caching for judgment-geometric verification.
\newblock Paper~38 in the \JuGeo{} series.

\bibitem{jugeo-paper65}
\JuGeo{} Development Team.
\newblock Integrating \JuGeo{} into CI/CD pipelines.
\newblock Paper~65 in the \JuGeo{} series.
\end{thebibliography}

\end{document}
